\documentclass[conference]{IEEEtran}

\usepackage{cite}
\usepackage{amsmath,amssymb,amsfonts}
\usepackage{algorithmic}
\usepackage{graphicx}
\usepackage{textcomp}
\usepackage{xcolor}
\usepackage{hyperref}

\begin{document}

\title{Paper Title}

\author{
    \IEEEauthorblockN{First Author}
    \IEEEauthorblockA{
        \textit{Department} \\
        \textit{University} \\
        City, Country \\
        email@example.com
    }
    \and
    \IEEEauthorblockN{Second Author}
    \IEEEauthorblockA{
        \textit{Department} \\
        \textit{University} \\
        City, Country \\
        email@example.com
    }
}

\maketitle

\begin{abstract}
    This document provides a template for IEEE conference papers.
    Replace this text with your abstract.
\end{abstract}

\begin{IEEEkeywords}
    keyword1, keyword2, keyword3
\end{IEEEkeywords}

\section{Introduction}

Your introduction here. Cite references like this~\cite{ref1}.

\section{Related Work}

Discuss related work here.

\section{Methodology}

Describe your approach.

\subsection{Problem Formulation}

Define the problem.

\subsection{Proposed Solution}

Detail your solution.

\section{Experiments}

\subsection{Experimental Setup}

Describe setup and datasets.

\subsection{Results}

Present your results. You can include figures:

% \begin{figure}[htbp]
%     \centering
%     \includegraphics[width=0.45\textwidth]{figures/example}
%     \caption{Example figure caption.}
%     \label{fig:example}
% \end{figure}

\subsection{Discussion}

Discuss the implications of your results.

\section{Conclusion}

Summarize your findings and future work.

\section*{Acknowledgment}

The authors would like to thank\ldots

\bibliographystyle{IEEEtran}
\bibliography{references}

\end{document}
